% Part: sets-functions-relations
% Chapter: sets
% Section: pairs-and-products

\documentclass[../../../include/open-logic-section]{subfiles}

\begin{document}

\olfileid{sfr}{set}{pai}
\olsection{جوڑے، ٹپل تے کارتیسی حاصلِ ضرب}

\begin{explain}
عنصراں نال تعیّن دے اصول توں نکلدا اے کہ سیٹ دے عنصراں دی کوئی ترتیب نہیں
ہوندی۔ سو، جے اسی ترتیب ظاہر کرنا چاہیئے، تاں \emph{ترتیب وار جوڑے}
$\tuple{x, y}$ ورتدے آں۔ بے ترتیب جوڑے $\{x, y\}$ وچ ترتیب دا کوئی فرق
نہیں پیندا: $\{x, y\} = \{y, x\}$۔ ترتیب وار جوڑے وچ فرق پیندا اے: جے
$x \neq y$ ہووے، تاں $\tuple{x, y} \neq \tuple{y, x}$ اے۔

سیٹاں دے نظریے وچ ترتیب وار جوڑیاں بارے ساڈا تصور کی ہونا چاہیدا اے؟ بنیادی
گل ایہ اے کہ اسی ایہ خیال قائم رکھنا چاہندے آں کہ ترتیب وار جوڑے اُدوں تے
صرف اُدوں اکو ہوندے نیں جدوں اوہناں دا پہلا عنصر اکو ہووے تے دوجا عنصر وی
اکو ہووے، یعنی:
\[
  \tuple{a, b}= \tuple{c, d}\text{ اُدوں تے صرف اُدوں، جدوں دونوں }a = c \text{ تے }b=d.
\]
سیٹاں دے نظریے وچ اسی وینر--کوراتوسکی دی تعریف نال ترتیب وار جوڑے دی
تعریف کر سکدے آں۔
\end{explain}

\begin{defn}[ترتیب وار جوڑا]\ollabel{wienerkuratowski}
	$\tuple{a, b} = \{\{a\}, \{a, b\}\}$۔
\end{defn}

\begin{prob}
	\olref[sfr][set][pai]{wienerkuratowski} ورت کے ثابت کرو کہ
	$\tuple{a, b}= \tuple{c, d}$ اُدوں تے صرف اُدوں ہوندا اے جدوں
	$a = c$ تے $b=d$، دونوں ہون۔
\end{prob}

\begin{explain}
ترتیب وار جوڑے دی تعریف طے کر کے اسی اوہنوں ہور سیٹاں دی تعریف لئی ورت
سکدے آں۔ مثال لئی، کدی کدی سانوں دو توں ودھ شےواں دے ترتیب وار سلسلے وی
چاہیدے ہوندے نیں، جیویں \emph{تِکے} $\tuple{x, y, z}$، \emph{چوکے}
$\tuple{x, y, z, u}$، تے ایس طرح اگے۔ تِکیاں نوں اسی خاص ترتیب وار جوڑے
سمجھ سکدے آں، جنہاں دا پہلا عنصر آپ اک ترتیب وار جوڑا اے:
$\tuple{x, y, z}$، $\tuple{\tuple{x, y},z}$ اے۔ چوکیاں لئی وی ایہو گل اے:
$\tuple{x,y,z,u}$، $\tuple{\tuple{\tuple{x,y},z},u}$ اے، تے ایس طرح اگے۔
عام طور تے اسی \emph{ترتیب وار $n$-ٹپل} $\tuple{x_1, \dots, x_n}$ دی گل
کردے آں۔

ترتیب وار جوڑیاں، یا ہور ترتیب وار $n$-ٹپلاں دے کجھ سیٹ ساڈے کم آؤن گے۔
\end{explain}

\begin{defn}[کارتیسی حاصلِ ضرب]
دتے ہوئے سیٹاں $A$ تے $B$ دا \emph{کارتیسی حاصلِ ضرب} $A \times B$ ایویں
متعین ہوندا اے:
\[
  A \times B = \Setabs{\tuple{x, y}}{x \in A \text{ تے } y \in B}.
\]
\end{defn}

\begin{ex}
جے $A = \{0, 1\}$ تے $B = \{1, a, b\}$ ہون، تاں اوہناں دا حاصلِ ضرب اے:
\[
A \times B = \{ \tuple{0, 1}, \tuple{0, a}, \tuple{0, b},
    \tuple{1, 1}, \tuple{1, a}, \tuple{1, b} \}.
\]
\end{ex}

\begin{ex}
جے $A$ اک سیٹ اے، تاں $A$ دا اپنے آپ نال حاصلِ ضرب، $A \times A$،
$A^2$ وی لکھیا جاندا اے۔ ایہ اوہناں \emph{سارے} جوڑیاں $\tuple{x, y}$ دا
سیٹ اے جنہاں لئی $x, y \in A$ اے۔ سارے تِکیاں $\tuple{x, y, z}$ دا
سیٹ $A^3$ اے، تے ایس طرح اگے۔ اسی اک بازگشتی تعریف دے سکدے آں:
\begin{align*}
  A^1 & = A\\
  A^{k+1} & = A^k \times A
\end{align*}
\end{ex}

\begin{prob}
سارے !!{element}s لکھو جو $\{1, 2, 3\}^3$ وچ نیں۔
\end{prob}

\begin{prop}\ollabel{cardnmprod}
جے $A$ دے $n$ !!{element}s تے $B$ دے $m$ !!{element}s ہون، تاں
$A \times B$ دے $n\cdot m$ عنصر ہون گے۔
\end{prop}

\begin{proof}
ہر !!{element}~$x$ لئی جو $A$ وچ اے، $m$ ایہو جہے !!{element}s نیں
جنہاں دی صورت $\tuple{x, y} \in A \times B$ اے۔ منّ لو کہ
$B_x = \Setabs{\tuple{x, y}}{y \in B}$۔ کیوں جو جدوں وی $x_1 \neq x_2$
ہووے، تاں $\tuple{x_1, y} \neq \tuple{x_2, y}$ اے، ایس لئی
$B_{x_1} \cap B_{x_2} = \emptyset$۔ پر جے $A = \{x_1, \dots, x_n\}$
ہووے، تاں $A \times B = B_{x_1} \cup \dots \cup B_{x_n}$ اے، تے ایس لئی
اوہدے $n\cdot m$ !!{element}s نیں۔

ایس گل نوں شکل وچ ویکھن لئی، \pnboblique{!!{element}s} نوں جو $A \times B$
وچ نیں، قطاراں تے کالماں دے جال وچ ترتیب نال رکھو:
\[
\begin{array}{rcccc}
  B_{x_1} = & \{\tuple{x_1, y_1} & \tuple{x_1, y_2} & \dots & \tuple{x_1, y_m}\}\\
  B_{x_2} = & \{\tuple{x_2, y_1} & \tuple{x_2, y_2} & \dots & \tuple{x_2, y_m}\}\\
  \vdots & & \vdots\\
  B_{x_n} = & \{\tuple{x_n, y_1} & \tuple{x_n, y_2} & \dots & \tuple{x_n, y_m}\}
\end{array}
\]
کیوں جو سارے $x_i$ وکھ وکھ نیں، تے سارے $y_j$ وی وکھ وکھ نیں، ایس جال
وچ کوئی دو جوڑے اکو نہیں، تے اوہناں دی گنتی $n\cdot m$ اے۔
\end{proof}

\begin{prob}
$k$ اُتے ریاضیاتی استقرا نال وکھاؤ کہ سارے $k \ge 1$ لئی، جے $A$ دے
$n$ !!{element}s ہون، تاں $A^k$ دے $n^k$ !!{element}s ہون گے۔
\end{prob}

\begin{ex}
جے $A$ اک سیٹ اے، تاں $A$ اُتے \emph{لفظ}، \pnboblique{!!{element}s} دا
کوئی وی سلسلہ اے جو $A$ توں لئے گئے ہون۔ سلسلے نوں اک $n$-ٹپل سمجھیا جا
سکدا اے، جہدے !!{element}s، $A$ توں نیں۔ مثال لئی، جے $A = \{a, b, c\}$
ہووے، تاں سلسلے ``$bac$'' نوں تِکّے~$\tuple{b, a, c}$ دے طور تے سمجھیا
جا سکدا اے۔ لفظ، یعنی علامتاں دے سلسلے، کمپیوٹر سائنس وچ بنیادی اہمیت
رکھدے نیں۔ رواج موجب، اسی \pnboblique{!!{element}s} نوں جو $A$ وچ نیں،
لمبائی~$1$ والے سلسلے گِندے آں، تے $\emptyset$ نوں لمبائی~$0$ والا سلسلہ۔
فیر $A$ اُتے \emph{سارے} لفظاں دا سیٹ اے:
\[
A^* = \{\emptyset\} \cup A \cup A^2 \cup A^3 \cup \dots
\]
\end{ex}

\end{document}
